% !PRINTED 137
\weberpage{137}
%<<CONT

By multiplication by $y^n$ and division by $a_n$, (21) then passes into
\begin{equation}
 y^n+\frac{a_{n-1}}{a_n}y^{n-1}+\cdots+\frac{a_1}{a_n}y+\frac{a_0}{a_n}=0,
\tag{23}
\end{equation}
and we can apply to the roots of this equation the theorem just stated, by assuming $a_0=0$, so that one of the roots of (23) vanishes. We thus obtain the theorem:

One can take $a_0$ in the equation (21) so small that one of its roots becomes large beyond all bounds, while the others differ arbitrarily little from the roots of the equation of degree $(n-1)$
\[
 a_1x^{n-1}+a_2x^{n-2}+\cdots+a_{n-1}x+a_n=0.
\]
This is also expressed by saying that, \gesperrt{as $a_0$ vanishes, one of the roots of (21) becomes infinite.}

If $a_1$ also vanishes, or $a_1$ and $a_2$, and so on, two or three roots become infinite.

This justifies the expression sometimes used, that an equation of degree $n$ passes, by the becoming-infinite of one root, into an equation of degree $(n-1)$.

\begin{center}
\rule{0.16\textwidth}{0.4pt}
\end{center}
